\documentclass[10pt,a4paper]{amsart}
\usepackage[margin=19mm]{geometry}
\usepackage{amsmath,amssymb,mathtools}
\usepackage[scr=rsfs,cal=euler]{mathalfa}
\usepackage{xfrac}
\usepackage[unicode,hidelinks,bookmarks=false]{hyperref}
\newcommand{\ie}{i.e.\ }
\newcommand{\cf}{cf.\ }
\newcommand{\resp}{respectively\ }
\newcommand{\Subsection}{\subsection}
\providecommand{\nobreakdash}{\nobreak\hbox{-}}
\setlength{\emergencystretch}{2em}
\multlinegap0pt
\usepackage{mathtools}
\usepackage{amssymb}
\usepackage[shortlabels]{enumitem}
\usepackage{bm}
\newtheorem{theorem}{Theorem}
\newtheorem{lemma}{Lemma}
\title{On Consecutive Non-primitive Elements over Finite Fields}
\author{Bidushi Sharma, DHIREN KUMAR BASNET*}
\date{Source: arXiv:2608.27318v1 (CC BY 4.0)}
\begin{document}
\maketitle

\noindent\emph{Source and licence.} On Consecutive Non-primitive Elements over Finite Fields. Version arXiv:2608.27318v1 (CC BY 4.0), distributed by arXiv under the Creative Commons Attribution 4.0 International licence, http://creativecommons.org/licenses/by/4.0/, as recorded in the arXiv licence metadata for this exact version and shown on the abstract page https://arxiv.org/abs/2608.27318v1. Full licence notice: \texttt{licenses/arxiv-2608.27318-CC-BY-4.0.txt}.\par\smallskip

\noindent\textit{Editorial scope.} Counted unit: the article's lemma L2.3 (source0.tex lines 269--271). The article's complete original proof of it is delivered with it (source0.tex lines 272--363, one \texttt{proof} environment of 92 lines). No mathematics is written, repaired or re-derived: the statement and the proof are the article's own lines, in the article's own notation, and no retained line is substituted or re-flowed.  Retained uncounted as the source-local context the proof needs: lines 167--172.  Nothing was omitted from any retained span, so the package carries no omission record.  No retained line contains a citation, so the wrapper carries no bibliography and no bibliography entry.  No retained line contains a figure environment, an image-inclusion command or drawing code, so no image is delivered and the package declares no asset dependency.  Editorial formatting only, in addition: an AMS \texttt{amsart} preamble that restates the article's own declarations and macros used by the retained lines, namely the environments \texttt{theorem}, \texttt{lemma} on separate counters, together with the standard packages those lines need.  The retained environments print in this document as Theorem 1, Lemma 1 in order of appearance, whereas the article prints the same items with the numbering of its own full document.  The complete native source of the article as distributed in the arXiv e-print for this exact version is bundled unchanged as \texttt{sources/208707/source0.tex}.
\begin{theorem}\label{T2.1}
Let $q$ be an even prime power and let $\ell$ be a prime divisor of $q-1$. Then $\mathbb{F}_q$ contains an $\mathrm{NT}\ell$ pair whenever
\[
q>9\ell^4.
\]
\end{theorem}

\begin{lemma}\label{L2.3}
Let $q=2^k$, where $k>3$ is an odd multiple of $3$. Then $\theta_q<\frac{6}{7}$, and $\mathbb{F}_q$ contains an $\mathrm{NT}7$ pair.
\end{lemma}

\begin{proof}
By hypothesis, $k=3(2m+1)$ for some $m> 1$. Since $k$ is an odd multiple of $3$, $7$ is the least prime divisor of $q-1$. Hence,
\[
\theta_q
\leq\frac{6}{7}.
\]
We claim that the inequality is strict. Suppose, to the contrary, that $\theta_q=\frac{6}{7}$ for some $q$.
Then $7$ must be the only prime divisor of $q-1$. Put $n=2m+1$, so that $n>1$ is odd. We have
\[
q-1=2^{3n}-1
=(2^n-1)(2^{2n}+2^n+1).
\]

Now suppose that $n$ is not divisible by $3$, then $2^n\not\equiv 1\pmod 7$, and thus $2^n-1$ has a prime divisor different from $7$, a contradiction.

Finally, suppose that $3\mid n$. Then 
\[
2^{2n}+2^n+1\equiv 1+1+1\equiv3\pmod 7,
\]
 which is again a contradiction.

Thus, $\theta_q<\frac{6}{7}$.
By Theorem~\ref{T2.1}, it remains only to verify the existence of an $\mathrm{NT}7$ pair for even prime powers satisfying $q\leq 9\cdot7^4=21609$. 

Under the present hypotheses, the only value that remains to be checked is $q=2^9$. Using SageMath we obtain the following data.
\begin{table}[h]
\centering
\scriptsize
\begin{tabular}{c|c|c|p{8cm}}
$q$ & $\#M_7(q)$ & $\theta_q$ & $\mathrm{NT}7$ pairs \\ \hline
$2^9$ & $18$ & $0.8454$ &
$\left(a^8+a^6+a^5+a^3+a,\,
a^8+a^6+a^5+a^3+a+1\right)$, \\
&&&
$\left(a^8+a^6+a^2+a+1,\,
a^8+a^6+a^2+a\right)$, \\
&&&
$\left(a^5+a^3+a,\,
a^5+a^3+a+1\right)$, \\
&&&
$\left(a^6+a^4+a^3+1,\,
a^6+a^4+a^3\right)$, \\
&&&
$\left(a^8+a^6+a^2+a,\,
a^8+a^6+a^2+a+1\right)$, \\
&&&
$\left(a^6+a^5+a^2+a,\,
a^6+a^5+a^2+a+1\right)$, \\
&&&
$\left(a^8+a^7+a^6+a^3+1,\,
a^8+a^7+a^6+a^3\right)$, \\
&&&
$\left(a^8+a^6+a^5+a^3+a+1,\,
a^8+a^6+a^5+a^3+a\right)$, \\
&&&
$\left(a^6+a^4+a^3,\,
a^6+a^4+a^3+1\right)$, \\
&&&
$\left(a^8+a^6+a^5+a^4+a,\,
a^8+a^6+a^5+a^4+a+1\right)$, \\
&&&
$\left(a^7+a^5+a^4+a^3+a^2+a,\,
a^7+a^5+a^4+a^3+a^2+a+1\right)$, \\
&&&
$\left(a^7+a^5+a^4+a^3+a^2+a+1,\,
a^7+a^5+a^4+a^3+a^2+a\right)$, \\
&&&
$\left(a^8+a^6+a^5+a^4+a+1,\,
a^8+a^6+a^5+a^4+a\right)$, \\
&&&
$\left(a^5+a^4+a^3+a^2,\,
a^5+a^4+a^3+a^2+1\right)$, \\
&&&
$\left(a^6+a^5+a^2+a+1,\,
a^6+a^5+a^2+a\right)$, \\
&&&
$\left(a^5+a^4+a^3+a^2+1,\,
a^5+a^4+a^3+a^2\right)$, \\
&&&
$\left(a^5+a^3+a+1,\,
a^5+a^3+a\right)$, \\
&&&
$\left(a^8+a^7+a^6+a^3,\,
a^8+a^7+a^6+a^3+1\right)$.
\end{tabular}
\caption{The $\mathrm{NT}7$ pairs in $\mathbb{F}_{2^9}$.}
\label{tab:nt7}
\end{table}


This completes our proof.
\end{proof}

\end{document}
